\chapter{Diseño y Desarrollo}
\label{cap:capitulo5}

Una vez definida la problemática abordada, los objetivos del proyecto y el marco tecnológico en el que se desarrolla, en este capítulo se describe el proceso seguido para construir la solución propuesta.
Se presentan los prototipos desarrollados hasta alcanzar la versión final funcional, describiendo su implementación y validación, así como los errores encontrados y las soluciones adoptadas.\\

Cabe resaltar que el código fuente oficial del \textit{Kernel} de \textit{Linux} se encuentra disponible en \url{https://www.kernel.org}, y que la versión del \textit{Kernel} utilizada en el proyecto es la \texttt{6.12.41-current-bcm2711}.
Esta precisión resulta importante, ya que el desarrollo y las referencias consultadas se basan en el código fuente de esta versión, pudiendo variar tanto la funcionalidad como la localización de algunos recursos respecto a otras versiones del \textit{Kernel}.\\

\section{Prototipos, Implementaciones y Validaciones}
\label{sec:prototipos-implementaciones-validaciones}

En esta sección se describen los prototipos desarrollados durante el proyecto desde una perspectiva conceptual, explicando los conocimientos teóricos aplicados en cada etapa y la forma en la que dichos conceptos se integran progresivamente en la solución final.
El objetivo no es reproducir aquí el código fuente, sino justificar el diseño seguido, detallar el funcionamiento de cada mecanismo y relacionar cada prototipo con la necesidad técnica que resuelve dentro del sistema de reporte seguro propuesto.\\

La implementación en detalle, los códigos y scripts utilizados junto con las salidas de validación generadas se recogen en el Apéndice~\ref{anexo:codigo-prototipos}.
Por otra parte, el código fuente completo del proyecto se encuentra también disponible en el repositorio público de \textit{GitHub} \url{https://github.com/nowelito28/TFG}.\\

\subsection{\textit{Hello World}}

\subsection{Fichero virtual \texttt{/proc/fdev}}

\subsection{\textit{File Handoff}}

\subsection{Mecanismo criptográfico de autenticación \textit{HMAC-SHA256}}

\subsection{Sistema de reporte seguro de los procesos en ejecución del sistema}

\section{Problemas y Errores encontrados}


\section{Objetivos cumplidos}
